% ============================================================
%  SECTION 3 — ALGORITHM PROJECT AND ANALYSIS
% ============================================================
\section{Algorithm Project and Analysis}
\label{sec:algorithms}

This section details the algorithmic design and formal asymptotic analysis of the four fractal dimension estimators. For each method, we first define a naive brute-force formulation, identify its computational bottleneck, and subsequently derive a structurally optimized counterpart. Throughout this analysis, $N = W \times H$ denotes the total image pixel count, $m = |\mathcal{S}|$ the number of evaluated scales, and $\varepsilon$ (or $L$) the current box side length. Asymptotic bounds follow standard Bachmann–Landau notation ($\mathcal{O}$, $\Theta$, $\Omega$).

% ------------------------------------------------------------
\subsection{Classical Box-Counting (BC)}
\label{subsec:bc}
% ------------------------------------------------------------

\subsubsection{Naive Formulation}

Algorithm~\ref{alg:bc_naive} implements Eq.~\eqref{eq:bc_count} through exhaustive cell scanning. For each scale $\varepsilon$, the grid contains $\Theta(N/\varepsilon^{2})$ cells. In the worst case (a fully foreground cell), evaluating a cell requires $\varepsilon^{2}$ pixel inspections; in the best case (first pixel is foreground), it requires one. The worst-case per-scale cost is:

\begin{equation}
  T_{\text{BC,naive}}(\varepsilon)
    = \Theta\!\left(\frac{N}{\varepsilon^{2}}\right)
      \cdot \mathcal{O}(\varepsilon^{2})
    = \mathcal{O}(N).
\end{equation}

The strict lower bound is $\Omega(N/\varepsilon^{2})$ since every cell must be visited. Summing over all $m$ scales, including the $\Theta(N)$ binarization pass, yields the total execution time:

\begin{equation}
  T_{\text{BC,naive}}(N,m)
    = \Theta(N) + \sum_{\varepsilon \in \mathcal{S}}
      \mathcal{O}(N)
    = \mathcal{O}(mN).
  \label{eq:bc_naive_time}
\end{equation}

Beyond the read-only input buffer, auxiliary space is $\Theta(1)$.

\begin{algorithm}[t]
\caption{Naive Classical Box-Counting}
\label{alg:bc_naive}
\begin{algorithmic}[1]
\Require Binary image $I[0\ldots W{-}1][0\ldots H{-}1]$, scale set $\mathcal{S}$
\Ensure Pairs $\{(\varepsilon,\, N(\varepsilon)) : \varepsilon \in \mathcal{S}\}$
\State Binarize $I$ via Otsu's threshold $t^{*}$ \cite{Otsu1979}
\For{each $\varepsilon \in \mathcal{S}$}
  \State $\mathit{count} \leftarrow 0$
  \For{$b_x \leftarrow 0$ \textbf{to} $\lceil W/\varepsilon \rceil - 1$}
    \For{$b_y \leftarrow 0$ \textbf{to} $\lceil H/\varepsilon \rceil - 1$}
      \State $\mathit{occupied} \leftarrow \textsc{False}$
      \For{$dx \leftarrow 0$ \textbf{to} $\varepsilon - 1$}
        \For{$dy \leftarrow 0$ \textbf{to} $\varepsilon - 1$}
          \State $px \leftarrow b_x \cdot \varepsilon + dx$;\; $py \leftarrow b_y \cdot \varepsilon + dy$
          \If{$px < W \;\wedge\; py < H \;\wedge\; I[px][py] = 1$}
            \State $\mathit{occupied} \leftarrow \textsc{True}$;\; \textbf{break}
          \EndIf
        \EndFor
      \EndFor
      \If{$\mathit{occupied}$}
        \State $\mathit{count} \leftarrow \mathit{count} + 1$
      \EndIf
    \EndFor
  \EndFor
  \State \textbf{yield} $(\varepsilon,\, \mathit{count})$
\EndFor
\State Estimate $D_B$ by least-squares regression on $\{(\log \varepsilon_i,\, \log N(\varepsilon_i))\}$
\end{algorithmic}
\end{algorithm}

% ............................................................
\subsubsection{Optimized Formulation}

The inner double loop of Algorithm~\ref{alg:bc_naive} represents the primary computational bottleneck. Algorithm~\ref{alg:bc_opt} eliminates this by pre-computing the integral image $P$ (Eq.~\eqref{eq:prefix}). Consequently, region sums are resolved in $\mathcal{O}(1)$ via Eq.~\eqref{eq:rect_sum}~\cite{Crow1984}, determining cell occupancy ($\sigma > 0$) without pixel scanning.

\begin{algorithm}[t]
\caption{Optimized Box-Counting (Integral Image)}
\label{alg:bc_opt}
\begin{algorithmic}[1]
\Require Binary image $I[0\ldots W{-}1][0\ldots H{-}1]$, scale set $\mathcal{S}$
\Ensure Pairs $\{(\varepsilon,\, N(\varepsilon)) : \varepsilon \in \mathcal{S}\}$
\State Binarize $I$ via Otsu's threshold $t^{*}$ \cite{Otsu1979}
\State Build integral image $P$ via Eq.~\eqref{eq:prefix} \hfill \Comment{$\Theta(N)$ time and space \cite{Crow1984}}
\For{each $\varepsilon \in \mathcal{S}$}
  \State $\mathit{count} \leftarrow 0$
  \For{$b_x \leftarrow 0$ \textbf{to} $\lceil W/\varepsilon \rceil - 1$}
    \For{$b_y \leftarrow 0$ \textbf{to} $\lceil H/\varepsilon \rceil - 1$}
      \State $x_1 \leftarrow b_x \cdot \varepsilon$;\; $y_1 \leftarrow b_y \cdot \varepsilon$
      \State $x_2 \leftarrow \min(x_1{+}\varepsilon{-}1, W{-}1)$;\; $y_2 \leftarrow \min(y_1{+}\varepsilon{-}1, H{-}1)$
      \If{$\sigma(x_1, y_1, x_2, y_2) > 0$} \hfill\Comment{Eq.~\eqref{eq:rect_sum}, $\mathcal{O}(1)$}
        \State $\mathit{count} \leftarrow \mathit{count} + 1$
      \EndIf
    \EndFor
  \EndFor
  \State \textbf{yield} $(\varepsilon,\, \mathit{count})$
\EndFor
\State Estimate $D_B$ by least-squares regression on $\{(\log \varepsilon_i,\, \log N(\varepsilon_i))\}$
\end{algorithmic}
\end{algorithm}

The prefix sum is computed once in $\Theta(N)$. Subsequently, traversing the $\Theta(N/\varepsilon^{2})$ cells per scale costs $\mathcal{O}(1)$ each:

\begin{equation}
  T_{\text{BC,opt}}(N, m)
    = \Theta(N)
    + \sum_{\varepsilon \in \mathcal{S}}
      \Theta\!\left(\frac{N}{\varepsilon^{2}}\right).
  \label{eq:bc_opt_time}
\end{equation}

For a geometric scale set (e.g., $\varepsilon \in \{2, 4, 8 \ldots\}$), the sum is dominated by the smallest scale, collapsing the total time complexity to $\Theta(N)$. This strictly improves upon the $\mathcal{O}(mN)$ naive bound. The required integral image introduces a $\Theta(N)$ memory footprint.

\begin{table}[h]
\centering
\caption{Asymptotic complexity of Classical Box-Counting.}
\label{tab:bc_complexity}
\begin{tabular}{lcc}
\hline
\textbf{Formulation} & \textbf{Time} & \textbf{Aux.\ Space} \\
\hline
Naive     & $\mathcal{O}(mN)$ & $\Theta(1)$ \\
Optimized & $\Theta(N)$       & $\Theta(N)$ \\
\hline
\end{tabular}
\end{table}

% ------------------------------------------------------------
\subsection{Gliding Box (GB)}
\label{subsec:gb}
% ------------------------------------------------------------

\subsubsection{Naive Formulation}

Algorithm~\ref{alg:gb_naive} implements the dense sliding window mass calculation defined in Eq.~\eqref{eq:gb_mass}. For a fixed scale $L$, scanning $L^2$ pixels across $\Theta(N)$ valid window positions yields a per-scale cost of:

\begin{equation}
  T_{\text{GB,naive}}(L)
    = \Theta(N) \cdot \Theta(L^{2})
    = \Theta(NL^{2}).
  \label{eq:gb_naive_scale}
\end{equation}

Summing over $m$ scales, the total execution time is heavily penalized by larger box sizes:

\begin{equation}
  T_{\text{GB,naive}}(N, m)
    = \Theta\!\left(N \sum_{L \in \mathcal{S}} L^{2}\right).
  \label{eq:gb_naive_total}
\end{equation}

The per-scale mass histogram occupies $\mathcal{O}(L_{\max}^{2})$ auxiliary space, which is generally negligible.

\begin{algorithm}[t]
\caption{Naive Gliding Box}
\label{alg:gb_naive}
\begin{algorithmic}[1]
\Require Binary image $I[0\ldots W{-}1][0\ldots H{-}1]$, scale set $\mathcal{S}$
\Ensure Pairs $\{(L,\, \mu(L),\, \Lambda(L)) : L \in \mathcal{S}\}$
\State Binarize $I$ via Otsu's threshold $t^{*}$ \cite{Otsu1979}
\For{each $L \in \mathcal{S}$}
  \State $H[0 \ldots L^{2}] \leftarrow 0$ \hfill\Comment{mass-frequency histogram}
  \For{$x \leftarrow 0$ \textbf{to} $W - L$}
    \For{$y \leftarrow 0$ \textbf{to} $H - L$}
      \State $\mathit{mass} \leftarrow 0$
      \For{$dx \leftarrow 0$ \textbf{to} $L - 1$}
        \For{$dy \leftarrow 0$ \textbf{to} $L - 1$}
          \State $\mathit{mass} \leftarrow \mathit{mass} + I[x{+}dx][y{+}dy]$
        \EndFor
      \EndFor
      \State $H[\mathit{mass}] \leftarrow H[\mathit{mass}] + 1$
    \EndFor
  \EndFor
  \State $T \leftarrow (W{-}L{+}1)(H{-}L{+}1)$
  \State Compute $\mu(L)$ and $\mu^{(2)}(L)$ from $H$ and $T$
  \State Compute $\Lambda(L)$ via Eq.~\eqref{eq:lacunarity}
  \State \textbf{yield} $(L,\, \mu(L),\, \Lambda(L))$
\EndFor
\State Estimate $D_B$ by least-squares regression on $\{(\log L_i,\, \log \mu(L_i))\}$
\end{algorithmic}
\end{algorithm}

% ............................................................
\subsubsection{Optimized Formulation}

Algorithm~\ref{alg:gb_opt} applies the integral-image optimization to the GB method. Pre-computing $P$ (Eq.~\eqref{eq:prefix}) reduces each dense window summation to $\mathcal{O}(1)$. If BC and GB operate in the same pipeline, this $\Theta(N)$ data structure is entirely shared.

\begin{algorithm}[t]
\caption{Optimized Gliding Box (Integral Image)}
\label{alg:gb_opt}
\begin{algorithmic}[1]
\Require Binary image $I[0\ldots W{-}1][0\ldots H{-}1]$, scale set $\mathcal{S}$, integral image $P$
\Ensure Pairs $\{(L,\, \mu(L),\, \Lambda(L)) : L \in \mathcal{S}\}$
\For{each $L \in \mathcal{S}$}
  \State $H[0 \ldots L^{2}] \leftarrow 0$
  \For{$x \leftarrow 0$ \textbf{to} $W - L$}
    \For{$y \leftarrow 0$ \textbf{to} $H - L$}
      \State $\mathit{mass} \leftarrow \sigma(x,\, y,\, x{+}L{-}1,\, y{+}L{-}1)$ \hfill\Comment{Eq.~\eqref{eq:rect_sum}, $\mathcal{O}(1)$}
      \State $H[\mathit{mass}] \leftarrow H[\mathit{mass}] + 1$
    \EndFor
  \EndFor
  \State $T \leftarrow (W{-}L{+}1)(H{-}L{+}1)$
  \State Compute $\mu(L)$ and $\mu^{(2)}(L)$ from $H$ and $T$
  \State Compute $\Lambda(L)$ via Eq.~\eqref{eq:lacunarity}
  \State \textbf{yield} $(L,\, \mu(L),\, \Lambda(L))$
\EndFor
\State Estimate $D_B$ by least-squares regression on $\{(\log L_i,\, \log \mu(L_i))\}$
\end{algorithmic}
\end{algorithm}

With $\mathcal{O}(1)$ mass evaluations, the per-scale complexity drops to the $\Theta(N)$ valid window positions:

\begin{equation}
  T_{\text{GB,opt}}(N, m)
    = \Theta(N) + \sum_{L \in \mathcal{S}} \Theta(N)
    = \Theta(mN).
  \label{eq:gb_opt_total}
\end{equation}

This represents a $\Theta(L^{2})$ acceleration per scale over the naive implementation. While BC visits $\Theta(N/\varepsilon^{2})$ disjoint cells, GB must visit $\Theta(N)$ overlapping positions regardless of scale. This algorithmic overhead enables GB to extract the richer mass-frequency distributions necessary for lacunarity estimation ($\Lambda$).

\begin{table}[h]
\centering
\caption{Asymptotic complexity of the Gliding Box method.}
\label{tab:gb_complexity}
\begin{tabular}{lcc}
\hline
\textbf{Formulation} & \textbf{Time} & \textbf{Aux.\ Space} \\
\hline
Naive     & $\Theta\!\left(N\!\sum_{L}L^{2}\right)$ & $\mathcal{O}(L_{\max}^{2})$ \\
Optimized & $\Theta(mN)$ & $\Theta(N)$ \\
\hline
\end{tabular}
\end{table}

\begin{table}[h]
\centering
\caption{BC vs.\ GB optimized formulations (per scale).}
\label{tab:bc_gb_comparison}
\resizebox{\columnwidth}{!}{%
\begin{tabular}{lccc}
\hline
\textbf{Method} & \textbf{Positions/scale} & \textbf{Cost/position} & \textbf{Per-scale cost} \\
\hline
BC (opt.) & $\Theta(N/\varepsilon^{2})$ & $\mathcal{O}(1)$ & $\Theta(N/\varepsilon^{2})$ \\
GB (opt.) & $\Theta(N)$ & $\mathcal{O}(1)$ & $\Theta(N)$ \\
\hline
\end{tabular}%
}
\end{table}

% ------------------------------------------------------------
\subsection{Differential Box-Counting (DBC)}
\label{subsec:dbc}
% ------------------------------------------------------------

\subsubsection{Naive Formulation}

Algorithm~\ref{alg:dbc_naive} directly implements the 3D grid partitioning detailed in Eq.~\eqref{eq:dbc_count}. Because it operates on grayscale data, no binarization is performed. Traversing the $\Theta(N/\varepsilon^2)$ cells requires an exhaustive $\Theta(\varepsilon^2)$ pixel scan per cell to determine the local intensity extrema $(I^{ij}_{\max}, I^{ij}_{\min})$.

\begin{algorithm}[t]
\caption{Naive Differential Box-Counting}
\label{alg:dbc_naive}
\begin{algorithmic}[1]
\Require Grayscale image $I[0\ldots W{-}1][0\ldots H{-}1]$, intensity range $G$, scale set $\mathcal{S}$
\Ensure Pairs $\{(\varepsilon,\, N(\varepsilon)) : \varepsilon \in \mathcal{S}\}$
\For{each $\varepsilon \in \mathcal{S}$}
  \State $\delta \leftarrow \lceil G \;/\; \lceil W/\varepsilon \rceil \rceil$
  \State $\mathit{total} \leftarrow 0$
  \For{$b_x \leftarrow 0$ \textbf{to} $\lceil W/\varepsilon \rceil - 1$}
    \For{$b_y \leftarrow 0$ \textbf{to} $\lceil H/\varepsilon \rceil - 1$}
      \State $I_{\max} \leftarrow -\infty$;\; $I_{\min} \leftarrow +\infty$
      \For{$dx \leftarrow 0$ \textbf{to} $\varepsilon - 1$}
        \For{$dy \leftarrow 0$ \textbf{to} $\varepsilon - 1$}
          \State $px \leftarrow b_x{\cdot}\varepsilon{+}dx$;\; $py \leftarrow b_y{\cdot}\varepsilon{+}dy$
          \If{$px < W \;\wedge\; py < H$}
            \State $I_{\max} \leftarrow \max(I_{\max},\, I[px][py])$
            \State $I_{\min} \leftarrow \min(I_{\min},\, I[px][py])$
          \EndIf
        \EndFor
      \EndFor
      \State $\mathit{total} \leftarrow \mathit{total} + \lceil (I_{\max}{-}I_{\min})/\delta \rceil + 1$
    \EndFor
  \EndFor
  \State \textbf{yield} $(\varepsilon,\, \mathit{total})$
\EndFor
\State Estimate $D_B$ by least-squares regression on $\{(\log \varepsilon_i,\, \log N(\varepsilon_i))\}$
\end{algorithmic}
\end{algorithm}

The resulting per-scale cost is $\Theta(N/\varepsilon^2) \cdot \Theta(\varepsilon^2) = \Theta(N)$, yielding a total time complexity of:

\begin{equation}
  T_{\text{DBC,naive}}(N, m) = \Theta(mN).
  \label{eq:dbc_naive_total}
\end{equation}

Although the asymptotic class mirrors naive BC, the hidden constant factor is strictly larger. Two intensity comparisons are required per pixel, and early-exit mechanisms are mathematically impossible since true extrema require scanning the entire cell. Auxiliary space remains $\Theta(1)$.

% ............................................................
\subsubsection{Optimized Formulation}

Algorithm~\ref{alg:dbc_opt} replaces the exhaustive $\Theta(\varepsilon^2)$ scan with offline 2D sparse tables~\cite{Liu2014}. The tables map regional minima and maxima queries to $\mathcal{O}(1)$ execution time.

\begin{algorithm}[t]
\caption{Optimized Differential Box-Counting (Sparse Tables)}
\label{alg:dbc_opt}
\begin{algorithmic}[1]
\Require Grayscale image $I[0\ldots W{-}1][0\ldots H{-}1]$, intensity range $G$, scale set $\mathcal{S}$
\Ensure Pairs $\{(\varepsilon,\, N(\varepsilon)) : \varepsilon \in \mathcal{S}\}$
\State Build 2D sparse table $\mathcal{T}_{\max}$ for range maximum \hfill \Comment{$\Theta(N \log N)$ time and space \cite{Liu2014}}
\State Build 2D sparse table $\mathcal{T}_{\min}$ for range minimum \hfill \Comment{$\Theta(N \log N)$ time and space}
\For{each $\varepsilon \in \mathcal{S}$}
  \State $\delta \leftarrow \lceil G \;/\; \lceil W/\varepsilon \rceil \rceil$
  \State $\mathit{total} \leftarrow 0$
  \For{$b_x \leftarrow 0$ \textbf{to} $\lceil W/\varepsilon \rceil - 1$}
    \For{$b_y \leftarrow 0$ \textbf{to} $\lceil H/\varepsilon \rceil - 1$}
      \State $x_1 \leftarrow b_x{\cdot}\varepsilon$;\; $y_1 \leftarrow b_y{\cdot}\varepsilon$
      \State $x_2 \leftarrow \min(x_1{+}\varepsilon{-}1,\,W{-}1)$;\; $y_2 \leftarrow \min(y_1{+}\varepsilon{-}1,\,H{-}1)$
      \State $I_{\max} \leftarrow \mathcal{T}_{\max}.\textsc{Query}(x_1, y_1, x_2, y_2)$ \hfill\Comment{$\mathcal{O}(1)$}
      \State $I_{\min} \leftarrow \mathcal{T}_{\min}.\textsc{Query}(x_1, y_1, x_2, y_2)$ \hfill\Comment{$\mathcal{O}(1)$}
      \State $\mathit{total} \leftarrow \mathit{total} + \lceil (I_{\max}{-}I_{\min})/\delta \rceil + 1$
    \EndFor
  \EndFor
  \State \textbf{yield} $(\varepsilon,\, \mathit{total})$
\EndFor
\State Estimate $D_B$ by least-squares regression on $\{(\log \varepsilon_i,\, \log N(\varepsilon_i))\}$
\end{algorithmic}
\end{algorithm}

Sparse table construction demands $\Theta(N \log N)$ time and space. Post-construction, traversing the $\Theta(N/\varepsilon^2)$ cells executes in unit time per cell. The total complexity is therefore bound by the preprocessing phase:

\begin{equation}
  T_{\text{DBC,opt}}(N, m)
    = \Theta(N \log N) + \sum_{\varepsilon \in \mathcal{S}} \Theta\!\left(\frac{N}{\varepsilon^{2}}\right)
    = \Theta(N \log N).
  \label{eq:dbc_opt_total}
\end{equation}

While computationally superior at large scales, DBC optimization introduces a severe $\Theta(N \log N)$ memory overhead, establishing a prominent time-memory tradeoff for grayscale analyses.

\begin{table}[h]
\centering
\caption{Asymptotic complexity of Differential Box-Counting.}
\label{tab:dbc_complexity}
\begin{tabular}{lcc}
\hline
\textbf{Formulation} & \textbf{Time} & \textbf{Aux.\ Space} \\
\hline
Naive     & $\Theta(mN)$        & $\Theta(1)$        \\
Optimized & $\Theta(N \log N)$  & $\Theta(N \log N)$ \\
\hline
\end{tabular}
\end{table}

% ------------------------------------------------------------
\subsection{Blanket Method (BM)}
\label{subsec:bm}
% ------------------------------------------------------------

\subsubsection{Naive Formulation}

Algorithm~\ref{alg:bm_naive} implements the sequential morphological erosion and dilation recurrences defined in Eqs.~\eqref{eq:blanket_u}--\eqref{eq:blanket_vol}. Each of the $\varepsilon_{\max}$ dilation levels requires examining the 4-connected neighborhood ($\mathcal{O}(1)$ work) for all $N$ pixels. 

\begin{algorithm}[t]
\caption{Naive Blanket Method}
\label{alg:bm_naive}
\begin{algorithmic}[1]
\Require Grayscale image $I[0\ldots W{-}1][0\ldots H{-}1]$, maximum dilation level $\varepsilon_{\max}$
\Ensure Pairs $\{(\varepsilon,\, A(\varepsilon)) : \varepsilon = 1,\ldots,\varepsilon_{\max}\}$
\State $u \leftarrow I$;\; $b \leftarrow I$ \hfill\Comment{initialize both blankets}
\State $V_{\text{prev}} \leftarrow \sum_{x,y} [u[x][y] - b[x][y]]$
\For{$\varepsilon \leftarrow 1$ \textbf{to} $\varepsilon_{\max}$}
  \State $u' \leftarrow$ new array of size $W \times H$
  \State $b' \leftarrow$ new array of size $W \times H$
  \For{$x \leftarrow 0$ \textbf{to} $W-1$}
    \For{$y \leftarrow 0$ \textbf{to} $H-1$}
      \State $u'[x][y] \leftarrow u[x][y] + 1$
      \State $b'[x][y] \leftarrow b[x][y] - 1$
      \For{each $(s,t) \in \mathcal{N}(x,y)$}
        \State $u'[x][y] \leftarrow \max(u'[x][y],\; u[s][t])$
        \State $b'[x][y] \leftarrow \min(b'[x][y],\; b[s][t])$
      \EndFor
    \EndFor
  \EndFor
  \State $u \leftarrow u'$;\; $b \leftarrow b'$
  \State $V \leftarrow \sum_{x,y}[u[x][y]-b[x][y]]$
  \State \textbf{yield} $\bigl(\varepsilon,\;(V - V_{\text{prev}})/2\bigr)$
  \State $V_{\text{prev}} \leftarrow V$
\EndFor
\State Estimate $D_B$ by least-squares regression on $\{(\log \varepsilon_i,\, \log A(\varepsilon_i))\}$
\end{algorithmic}
\end{algorithm}

The total execution time evaluates to:

\begin{equation}
  T_{\text{BM,naive}}(N, \varepsilon_{\max}) = \Theta(N \cdot \varepsilon_{\max}).
  \label{eq:bm_naive_total}
\end{equation}

Because blanket surfaces strictly depend on their prior states, parallelizing across scales is impossible. Auxiliary space requires $\Theta(N)$ to maintain the active and previous state buffers.

% ............................................................
\subsubsection{Optimized Formulation}

Evaluating morphological extrema across an expanded $(2\varepsilon+1) \times (2\varepsilon+1)$ window typically incurs an $\Theta(\varepsilon^2)$ penalty per pixel. Algorithm~\ref{alg:bm_opt} sidesteps this by exploiting the \textit{separability} of maximum and minimum filters. A 2D extrema window decomposes into successive 1D row and column passes. Using Lemire's monotonic deques \cite{Lemire2006}, each 1D pass resolves in $\Theta(N)$ amortized time.

\begin{algorithm}[t]
\caption{Optimized Blanket Method (Separable Monotonic Deque)}
\label{alg:bm_opt}
\begin{algorithmic}[1]
\Require Grayscale image $I[0\ldots W{-}1][0\ldots H{-}1]$, maximum dilation level $\varepsilon_{\max}$
\Ensure Pairs $\{(\varepsilon,\, A(\varepsilon)) : \varepsilon = 1,\ldots,\varepsilon_{\max}\}$
\State $u \leftarrow I$;\; $b \leftarrow I$
\State $V_{\text{prev}} \leftarrow \sum_{x,y} [u[x][y] - b[x][y]]$
\For{$\varepsilon \leftarrow 1$ \textbf{to} $\varepsilon_{\max}$}
  \State $u_{\text{inc}} \leftarrow u + 1$ \hfill\Comment{element-wise increment, $\Theta(N)$}
  \State $b_{\text{dec}} \leftarrow b - 1$
  \State $u' \leftarrow \textsc{SlidingMax2D}(u,\, K{=}3)$ \hfill\Comment{separable deque, $\Theta(N)$ \cite{Lemire2006}}
  \State $b' \leftarrow \textsc{SlidingMin2D}(b,\, K{=}3)$ \hfill\Comment{separable deque, $\Theta(N)$}
  \State $u \leftarrow \max(u_{\text{inc}},\, u')$ \hfill\Comment{element-wise, $\Theta(N)$}
  \State $b \leftarrow \min(b_{\text{dec}},\, b')$
  \State $V \leftarrow \sum_{x,y} [u[x][y] - b[x][y]]$
  \State \textbf{yield} $\bigl(\varepsilon,\; (V - V_{\text{prev}})/2\bigr)$
  \State $V_{\text{prev}} \leftarrow V$
\EndFor
\State Estimate $D_B$ by least-squares regression on $\{(\log \varepsilon_i,\, \log A(\varepsilon_i))\}$
\end{algorithmic}
\end{algorithm}

The total asymptotic time remains equivalent to the naive class:

\begin{equation}
  T_{\text{BM,opt}}(N, \varepsilon_{\max}) = \Theta(N \cdot \varepsilon_{\max}).
  \label{eq:bm_opt_total}
\end{equation}

However, the separable formulation guarantees that the computational constant remains flat regardless of the structuring element size. Table~\ref{tab:all_complexity} consolidates the bounds across all evaluated methods.

\begin{table}[h]
\centering
\caption{Asymptotic complexity of the Blanket Method.}
\label{tab:bm_complexity}
\begin{tabular}{lcc}
\hline
\textbf{Formulation} & \textbf{Time} & \textbf{Aux.\ Space} \\
\hline
Naive     & $\Theta(N\varepsilon_{\max})$ & $\Theta(N)$ \\
Optimized & $\Theta(N\varepsilon_{\max})$ & $\Theta(N)$ \\
\hline
\end{tabular}
\end{table}

\begin{table}[h]
\centering
\caption{Consolidated asymptotic complexity of all four methods (optimized formulations).}
\label{tab:all_complexity}
\begin{tabular}{lccc}
\hline
\textbf{Method} & \textbf{Input} & \textbf{Time} & \textbf{Aux.\ Space} \\
\hline
BC  (opt.) & Binary    & $\Theta(N)$           & $\Theta(N)$        \\
GB  (opt.) & Binary    & $\Theta(mN)$          & $\Theta(N)$        \\
DBC (opt.) & Grayscale & $\Theta(N \log N)$    & $\Theta(N \log N)$ \\
BM  (opt.) & Grayscale & $\Theta(N\varepsilon_{\max})$ & $\Theta(N)$ \\
\hline
\end{tabular}
\end{table}